\documentclass[12pt]{article}

\input{../../_assets/preambulo.tex}


\begin{document}

    % 1. Foto de fondo
    % 2. Título
    % 3. Encabezado Izquierdo
    % 4. Color de fondo
    % 5. Coord x del titulo
    % 6. Coord y del titulo
    % 7. Fecha

    
    \input{../../_assets/portada}
    \portadaExamen{ffccA4.jpg}{Ecuaciones\\Diferenciales II\\Examen XX}{Ecuaciones Diferenciales II. Examen XX}{MidnightBlue}{-8}{28}{2026}{}

    \begin{description}
        \item[Asignatura] Ecuaciones Diferenciales II.
        \item[Curso Académico] 2021/2022.
        \item[Grado] Doble Grado en Ingeniería Informática y Matemáticas.
        \item[Grupo] Único.
        \item[Profesor] Rafael Ortega Ríos.
        \item[Descripción] Convocatoria Extraordinaria.
        \item[Fecha] 6 de julio de 2022.
        % \item[Duración] ---.
    
    \end{description}
    \newpage


    % ------------------------------------
    
    \begin{ejercicio}
        Construye una sucesión de funciones $\{f_n\}$, $f_n : \left[0, 1\right] \to \bb{R}$, que converja puntualmente a la función $f \equiv 0$ pero no lo haga de manera uniforme.
    \end{ejercicio}

    \begin{ejercicio}
        Resuelve el problema de valores iniciales
        \[
        x' = 1 + x^2, \quad x(0) = x_0,
        \]
        con $x_0 \in \bb{R}$. Determina en cada caso el intervalo maximal.
    \end{ejercicio}

    \begin{ejercicio}
        Se considera el problema de valores iniciales
        \[
        x'' + \sqrt{1 + x^2} = 0, \quad x(0) = 3, \quad x'(0) = 0.
        \]
        Demuestra que hay una única solución maximal que está definida en todo $\bb{R}$.
    \end{ejercicio}

    \begin{ejercicio}
        Encuentra todas las funciones continuas $x : \bb{R} \to \bb{R}$ que cumplen la ecuación integral
        \[
        x(t) = (\cos t) \int_0^t x(s) \, ds, \quad t \in \bb{R}.
        \]
    \end{ejercicio}

    \begin{ejercicio}~
        \begin{enumerate}[label=\alph*)]
            \item Prueba la desigualdad $\sqrt{a} \leq a + 1$ para cada $a \in \left[0, \infty\right[$.
            \item Se considera una sucesión de funciones continuas $x_n : \left[0, 1\right] \to \bb{R}$, que cumplen, para cada $n = 1, 2, \dots$,
            \[
            x_n(t) = \frac{1}{n} + \int_0^t |x_n(s)|^{\nicefrac{1}{2}} \, ds.
            \]
            Demuestra que la sucesión $\{x_n\}$ es uniformemente acotada y equicontinua.
        \end{enumerate}
    \end{ejercicio}

    \begin{ejercicio}
        La solución del problema
        \[
        x' = 4x(1 - x - \veps y), \quad y' = 6y(2 + \veps x - y), \quad x(0) = \veps + 1, \quad y(0) = 2\veps + 2
        \]
        se denota por $(x(t, \veps), y(t, \veps))$.
        \begin{enumerate}[label=\alph*)]
            \item Demuestra que si $|\veps|$ es lo bastante pequeño entonces esta solución está definida en $\left[-1, 1\right]$.
            \item Calcula, si existe, $\lim\limits_{\veps \to 0} x(t, \veps)$.
            \item Calcula, si existe, $\displaystyle \del{x}{\veps}\left(\frac{1}{2}, 0\right)$.
        \end{enumerate}
    \end{ejercicio}

    \newpage
    \setcounter{ejercicio}{0} % Reiniciar contador de ejercicios
    % \noindent
    % \textbf{Solución.}

\end{document}